\begin{abstract}

This thesis examines cybersecurity modelling techniques that are iterative, interdisciplinary, and incorporate different modelling and analysis tools. We developed a developed a middleware \textit{CyberVITBridge} and scaffold a feedback-loop workflow between a simulation software CyberBattleSim and the model-checking tool VITAMIN. This demonstrates a novel method for evaluating active defensive rather than a static defensive strategies in CyberBattleSim. However, these models lack compositionality and compatibility. Using compositional game theory (CGT) to build another class of cybersecurity models based on game theory, we demonstrate how these models can be composed and extended in an iterative way. After reviewing the current implementation of CGT as a Haskell library, we construct novel design patterns to streamline the construction of these models. Driven by software-design methodologies, we demonstrate how to rapidly prototype models-as-code and present a methodology for collaboratively maintaining and developing a library of models using a test-driven development cycle.  

\end{abstract}


